\chapter{Installation}
\label{ch:install}

\section{From PyPI}

The usual way:

\begin{lstlisting}
pip install poraque
\end{lstlisting}

\section{From source}

What you want if you intend to change anything --- the editable install means
an edit to the source takes effect without reinstalling:

\begin{lstlisting}
git clone https://github.com/seixas-research/poraque.git
cd poraque
pip install -e .
\end{lstlisting}

Either way this installs NumPy, SciPy, ASE, Matplotlib, PyYAML and PyTorch.
Symbolic distillation (Chapter~\ref{ch:symbolic}) needs an optional extra,
which pulls in PySR and its Julia toolchain:

\begin{lstlisting}
pip install "poraque[symbolic]"
\end{lstlisting}

\section{Console commands}
\label{sec:console-scripts}

Installing also registers five commands, which run from \emph{any} directory
once the environment is active --- no path to a script, no \opt{cd} into the
repository:

\begin{center}
\begin{tabular}{@{}lp{9.2cm}@{}}
\toprule
Command & Does \\
\midrule
\opt{poraque-train} & trains the operators \\
\opt{poraque-inference} & predicts a new structure \\
\opt{poraque-mp} & downloads densities from the Materials Project \\
\opt{poraque-committee} & \textbf{calibration}: does ensemble disagreement
predict error? \\
\opt{poraque-active-learning} & \textbf{selection}: which structures should
the next DFT runs be spent on? \\
\opt{poraque-vasp} & writes the VASP inputs that read a prediction back:
\opt{bands}, \opt{dos}, \opt{energy}; and \opt{chgcar}, which writes a field
store out as a \file{CHGCAR} \\
\bottomrule
\end{tabular}
\end{center}

\begin{pnote}[title={The last two are not the same command twice}]
They are the two halves of one loop, and they differ in the data they run on
--- which decides everything else.

\opt{poraque-committee} runs on a \textbf{labelled} dataset (inputs
\emph{and} targets), so the true error is available to correlate the
disagreement against. It answers \emph{is this measure worth trusting?} and
produces a Spearman coefficient. It costs nothing: the DFT is already done.

\opt{poraque-active-learning} runs on an \textbf{unlabelled} pool (inputs
only, targets not yet computed). It answers \emph{which structures do I
compute next?} and produces a ranking and a transfer into the training set.
It costs the DFT runs it selects.

Run the first one first. Both share the same committee, the same divergence
and the same ranking table, which is why their output looks alike --- what
differs is what the number is \emph{for}.
\end{pnote}

Each is the \opt{main()} of the script of the same name, with the hyphen
written as an underscore: \opt{python scripts/poraque\_train.py} is equivalent
to \opt{poraque-train} and needs nothing installed. Both forms take identical
arguments and behave identically.

\begin{pnote}
Relative paths inside a configuration file (\file{data/vasp}, \file{models/})
are resolved against the \textbf{working directory}, not the installation. Run
these commands from the project root, or give absolute paths.
\end{pnote}

If a command is not found after an upgrade, re-run \opt{pip install -e .}:
entry points are registered when the package is installed, not when it is
imported.

\section{Python version}

\begin{pnote}
Poraqu\^e requires \textbf{Python 3.11 or newer}. Every module is verified
against the 3.11 grammar. There is deliberately no upper bound, so newer
interpreters are supported.
\end{pnote}

\section{Checking the installation}

\begin{lstlisting}
pytest
python -m poraque.ml.device --check
python -c "from poraque.ml.device import available_devices; print(available_devices())"
\end{lstlisting}

The device list is ordered by preference, so the first entry is what
\opt{device: auto} will select --- \opt{cuda}, then \opt{mps}, then \opt{cpu}.
\opt{pytest -m "not gpu"} skips the tests that need an accelerator, rather than
reporting them as passes on a machine that has none.

\section{NVIDIA GPUs}
\label{sec:cuda}

\file{pyproject.toml} asks for \opt{torch>=2.0} and deliberately says nothing
about CUDA: which build is right depends on the machine, and pinning one would
break Apple Silicon, which wants the newest torch. But the choice has to be
made somewhere, because \textbf{the wrong one does not raise --- it gives a
silent CPU run}.

Two things have to line up: \textbf{the driver}, which must be new enough for
the wheel's CUDA runtime (\opt{nvidia-smi} prints the highest CUDA version it
supports, top right), and \textbf{the GPU architecture}, which must appear in
\opt{torch.cuda.get\_arch\_list()}.

\begin{center}
\begin{tabular}{@{}lll@{}}
\toprule
GPU & capability & wheel \\
\midrule
P100 & sm\_60 & \opt{cu118} / \opt{cu121} \\
\textbf{V100} & \textbf{sm\_70} & \textbf{\opt{cu126}} (CUDA 13 dropped Volta) \\
T4, RTX 20xx & sm\_75 & \opt{cu126} or newer \\
A100 & sm\_80 & any wheel $\geq$ \opt{cu118} \\
RTX 30xx & sm\_86 & any wheel $\geq$ \opt{cu118} \\
H100 & sm\_90 & \opt{cu121} or newer \\
B200, RTX 50xx & sm\_100 / sm\_120 & \opt{cu128} or newer \\
\bottomrule
\end{tabular}
\end{center}

\begin{lstlisting}
pip install torch --index-url https://download.pytorch.org/whl/cu126
pip install poraque
\end{lstlisting}

The second row is not hypothetical. \opt{pip install poraque} into a clean
environment on Santos Dumont brought \opt{torch 2.13.0+cu130}, which failed
twice over: the driver (560.35.03, CUDA 12.6) was too old for a \opt{cu130}
runtime, and \opt{sm\_70} is not in that build's \opt{arch\_list} at all. Even
with a newer driver the V100 would have had no kernel to run --- \emph{no
kernel image is available for execution on the device}, at the first forward
pass, after the queue time had been spent.

\subsection{Check before submitting to a queue}

\begin{lstlisting}
python -m poraque.ml.device --check || exit 1
\end{lstlisting}

It prints the torch build, its CUDA runtime, where it was imported from, the
architectures it carries kernels for, and one line per GPU saying whether this
build can use it --- then exits non-zero if the requested device is not usable.
Three seconds, against a slot in the GPU queue.

The same check belongs in the configuration, where it stops the run rather than
the script:

\begin{lstlisting}
training:
  device: cuda
  strict_device: true    # abort instead of falling back to the CPU
\end{lstlisting}

\begin{pnote}
Without it, a job that cannot reach its GPU takes its place in the queue, warns
into a log nobody reads until afterwards, and trains on the CPU \emph{inside}
the GPU allocation until the wall clock ends. \opt{strict\_device} is off by
default because a workstation wants the fallback --- a config written on a
machine with a GPU should still run on a laptop --- and should be on in every
batch job.
\end{pnote}

\subsection{\file{\textasciitilde/.local} outranks the active environment}

A \opt{torch} installed in \file{\textasciitilde/.local/lib/pythonX.Y/site-packages}
\textbf{wins} over the one in an active conda environment or venv. On a shared
cluster that is easy to acquire by accident and hard to notice: the environment
a job activated is not necessarily the one it ran.

\begin{lstlisting}
export PYTHONNOUSERSITE=1
\end{lstlisting}

Poraqu\^e's start-up banner prints torch's version, its CUDA build \emph{and its
install directory} for exactly this reason, so a mismatch appears in the run's
own log instead of becoming a four-hour CPU run.

\subsection{One GPU}

Poraqu\^e trains on a single device; request one. There is no
\opt{DistributedDataParallel} path, so a job allocated four GPUs will see four
and use \opt{cuda:0}. That is deliberate: for fields of this size the optimiser
step is tens of milliseconds on a model of a few megabytes, and the constraint
is getting data to it rather than the arithmetic. See
\opt{data.cache\_in\_memory} in the Configuration chapter, which is worth more
here than a second GPU would be.

\section{Faster CPU inference}
\label{sec:cbackend}

Poraqu\^e ships a small C kernel for the spectral contraction --- the one part
of a Fourier layer that PyTorch runs poorly at batch~1, which is the shape
every prediction has. It is \textbf{optional}: without it everything works and
falls back to \opt{torch.einsum}.

There is nothing to configure. The kernel is compiled on first use --- about
half a second, once --- and cached in \file{\textasciitilde/.cache/poraque}.
To do that compile now, and check it:

\begin{lstlisting}
python -m poraque.ml.backend --benchmark
\end{lstlisting}

\begin{lstlisting}
  C spectral backend: loaded from ...poraque_spectral_89feecc6.dylib (pthreads)
  agreement with torch.einsum: 2.76e-07 relative (float32 rounding is ~1e-7)

  one contraction (batch 1, width 32, modes 12^3):
    torch.einsum (4 threads)     4.528 ms
    C, serial                    0.606 ms     7.5x
    C, 4 pthreads                0.200 ms    22.6x
\end{lstlisting}

All it needs is a C compiler: \opt{xcode-select --install} on macOS, or
\opt{build-essential} on Debian/Ubuntu. Add \opt{--rebuild} after editing the
kernel.

\begin{center}
\begin{tabular}{@{}lp{9.4cm}@{}}
\toprule
Whole-model inference & 2--3.4$\times$ faster at $24^3$--$32^3$, tapering to
$\sim$1$\times$ at $96^3$ where the FFTs dominate \\
Training & unaffected: the kernel records nothing on the autograd tape, so it
is used only under \opt{torch.no\_grad()} \\
Accuracy & identical to float32 rounding ($\sim10^{-7}$ relative), checked on
every call path by the test-suite \\
To disable & \opt{PORAQUE\_C\_BACKEND=0} \\
\bottomrule
\end{tabular}
\end{center}

\opt{poraque-inference} prints which path it used, so a run that silently fell
back is visible rather than merely slow.

\section{Optional components}

\begin{center}
\begin{tabular}{@{}lll@{}}
\toprule
Component & Needed for & Install \\
\midrule
CUDA build of PyTorch & NVIDIA GPUs & \S\ref{sec:cuda} \\
\opt{pdflatex} / \opt{latexmk} & automatic PDF reports & TeX Live or MacTeX \\
\bottomrule
\end{tabular}
\end{center}

Apple Silicon needs nothing extra: the Metal backend ships with the standard
PyTorch wheel and is selected automatically.
